\section{Design}
\label{sec:design}

We propose \sys, a minimal DSN design that instantiates the data placement mechanism introduced in \cref{sec:approach}.
We commonly assume that there exists a kind of token that is economically favored by the rational nodes.
For example, the storage cost can be overly covered by the token rewarded for storage, and so on.

\subsection{On-chain Control Plane}

The smart contract of \sys maintains two on-chain states: current staking nodes and current storing data hashes.
The former state determines the number of nodes, which is used to determine current sample rate $p$.
The smart contract accepts three kinds of transactions: staking transactions for node joining, storing transaction for client storing data\footnotemark, and proving transaction for nodes to prove storage and gain rewards.
\footnotetext{For simplicity, we omit data expiration and removal.}

The staking transaction add nodes to the current staking nodes.
The storing transaction add data hash to the current storing data hashes.
On every block height, a random (node, data) pair is sampled from the states with the block hash as the seed, so the sampled pair cannot be predicted.
If the sampled node is endorsed by the data's placement group, the node can submit a proving transaction with VRF proof and the stored fragment within a time window.
The smart contract validates the node and the data hash in the transaction, then runs $\VRFVerify$ and $\ECVerify$ for validation as discussed in \cref{sec:approach:ec}.
If the validation passes, the submitting node is rewarded with token.

\subsection{Client Operations}

To store a data block, client submits a storing transaction with data hash and the token paid for storage.
When the nodes execute the storing transaction, they run $\VRFGen$ locally with their $\sk$ and current sample rate $p$ to determine whether they are endorsed.
If so, they query the client with VRF proofs.
The client first run $\VRFVerify$ to validate the nodes, and then encodes the fragments according to the proof hashes as indices and responds to the nodes.
The client can conclude the store operation after sufficient number of encoded fragments have been sent.

To retrieve a data block, client sends a gossip message to \sys to look for encoded fragments.
The nodes who store the fragments send them to the client.
The client then validates the fragments similar to the proving transaction and decodes the data block with the validated ones.
In this work, we omit the incentive during data retrieval.

\subsection{Node Bootstrapping}

Repairing for redundancy loss usually happens naturally when new nodes bootstrap into \sys.
After synchronizing the smart contract state, the nodes run $\VRFGen$ against the current storing data hashes.
Then, for each placement group it is endorsed, it queries the encoded fragments and decodes the data block similar to a client retrieval operation.
Then it encodes the fragment it should store according to the VRF proof hash and stores it locally.

\subsection{Configure Sample Rate}

The sample rate $p$ is periodically reconfigured based on Byzantine faulty rate $f$, maximum simultaneous benign failure size in a placement group $f'$, EC recover threshold $k$ and the number of nodes $N$.
All $f$, $f'$ and $k$ are preconfigured system parameters, and $N$ is derived from the current staking nodes state.
The formula of $p$ is shown in \XXX.
If the $p$ decreases because of $N$ grows, every node rerun $\VRFGen$ on the data of stored fragments and discard the fragments if the node is no longer endorsed.
Otherwise, if the $p$ increases due to the rare case of system shrinking, every node rerun the bootstrap routine and store additional fragments.
According to the random sampling definition (\cref{sec:approach:sample}), nodes only need to discard fragments when $p$ decreases and only need to store more fragments when $p$ increases.
This design minimizes the number of transferred encoded fragments during reconfiguration.